% Section content template for individual Educative lessons
% This file contains the actual content for a specific section
% Generated from Educative API section content

% Section: Code for the Online Stock Brokerage System
% Section ID: 5620382120017920
% Section Slug: code-for-the-online-stock-brokerage-system
% Generated: 2025-12-14T12:54:37.767151

We've reviewed different aspects of the online stock brokerage system and observed the attributes attached to the problem using various UML diagrams. Let 's explore the more practical side of things, where we will work on implementing the online stock brokerage system using multiple languages. This is usually the last step in an object-oriented design interview process.
\\
We have chosen the following languages to write the skeleton code of the different classes present in the online stock brokerage system:

\begin{itemize}
\item
\protect\phantomsection\label{hK8rctAX5bkWJx4CsuxKu}
Java
\item
\protect\phantomsection\label{8-N6jGwaULXxsFZOq1Gdl}
C\#
\item
\protect\phantomsection\label{LRHsGb6DrwjRFCF60Sw7z}
Python
\item
\protect\phantomsection\label{3-08WR2ZhsdxgOwHuHWiM}
C++
\item
\protect\phantomsection\label{4fzmPg7_DQSw90tlPrno0}
JavaScript
\end{itemize}
\\
If you want to skip directly to the implementation and see the complete workflow, simply scroll down or click~\hyperref[Executable-code-Online-stock-brokerage-system]{Executable Code: Online Stock Brokerage System}.

\subsection{Online stock brokerage system classes}\label{iB16EEgsaXif-uU-HRj7B}
\\
This section will provide the skeleton code of the classes designed in the class diagram lesson.
\begin{quote}
\textbf{Note:} For simplicity, we are not defining getter and setter functions. The reader can assume that all class attributes are private, accessed through their respective public getter methods, and modified only through their public method functions.
\end{quote}
\subsubsection{Enumerations and custom data type}\label{rOVMnzwZwv3Sg9LbBzNXQ}
\\
First, we will define all the enumerations required in the stock brokerage system. According to the class diagram, four enumerations are used in the system, i.e., \texttt{OrderStatus}, \texttt{TimeEnforcementType}, \texttt{AccountStatus},and \texttt{ErrorCode}\emph{.} The code to implement these enumerations and custom data types is as follows:
\begin{quote}
\textbf{Note:} JavaScript does not support enumerations, so we will use the \texttt{Object.freeze()} method as an alternative that freezes an object and prevents further modifications.
\end{quote}

\textbf{Definition of enums and custom data types}

\begin{lstlisting}[language=Java]
// Enumeration
enum OrderStatus {
  OPEN, 
  FILLED, 
  PARTIALLY_FILLED, 
  CANCELED
}

enum TimeEnforcementType {
  GOOD_TILL_CANCELED, 
  FILL_OR_TERMINATE, 
  IMMEDIATE_OR_CANCEL, 
  ON_THE_OPEN,
  ON_THE_CLOSE
}

enum AccountStatus {
  ACTIVE, 
  CLOSED, 
  CANCELED, 
  BLACKLISTED, 
  NONE
}

public enum ErrorCode {
    SUCCESS,
    INSUFFICIENT_FUNDS,
    INVALID_STOCK,
    ORDER_REJECTED
}

// Custom Address data type class
public class Address {
  private int zipCode;
  private String address;
  private String city;
  private String state;
  private String country;
}
\end{lstlisting}

\subsubsection{Account}\label{J-I53vTuIEEZkuWsUQZsK}
\\
The \texttt{Account} class will be an abstract class with the actors, \texttt{Admin} and \texttt{Member}, as child classes. The definition of these classes is given below:

\textbf{Account and its child classes}

\begin{lstlisting}[language=Java]
// Enumeration
enum OrderStatus {
  OPEN, 
  FILLED, 
  PARTIALLY_FILLED, 
  CANCELED
}

enum TimeEnforcementType {
  GOOD_TILL_CANCELED, 
  FILL_OR_TERMINATE, 
  IMMEDIATE_OR_CANCEL, 
  ON_THE_OPEN,
  ON_THE_CLOSE
}

enum AccountStatus {
  ACTIVE, 
  CLOSED, 
  CANCELED, 
  BLACKLISTED, 
  NONE
}

public enum ErrorCode {
    SUCCESS,
    INSUFFICIENT_FUNDS,
    INVALID_STOCK,
    ORDER_REJECTED
}

// Custom Address data type class
public class Address {
  private int zipCode;
  private String address;
  private String city;
  private String state;
  private String country;
}
\end{lstlisting}

\subsubsection{Watch list and stock}\label{LskWRVeFX0nb_9XnP-vEI}
\\
A \texttt{Watchlist} class is a list of stocks that an investor~ monitors to profit from price drops. The
\texttt{Stock} class is an equity or security that represents a portion of the issuing company's ownership.

\textbf{The Watchlist and the Stock classes}

\begin{lstlisting}[language=Java]
public class Watchlist {
  private String name;
  private List<Stock> stocks;
}

public class Stock {
  private String symbol;
  private double price;
}
\end{lstlisting}

\subsubsection{Search and stock inventory}\label{LRBfPcbLUEgDNNOwVnyuQ}
\\
The \texttt{StockInventory} class implements the \texttt{Search} interface.

\textbf{The Search interface and the StockInventory class}

\begin{lstlisting}[language=Java]
public interface Search {
    // Interface method (does not have a body)
    public Stock searchSymbol(String symbol);
}

public class StockInventory implements Search {
  private String inventoryName;
  private Date lastUpdate;
  public Stock searchSymbol(String symbol);
  public boolean deductStock(String symbol, double quantity);
  public boolean sendOrderDetails(Order order);
}
\end{lstlisting}

\subsubsection{Stock position and stock lot}\label{Hpa-z3TKp0HOtc1gDlebp}
\\
All the stocks the user owns will be included in the \texttt{StockPosition} class. A member may purchase various lots of the same stock on various dates. The
\texttt{StockLot} class will represent these particular lots.

\textbf{The StockPosition and StockLot classes}

\begin{lstlisting}[language=Java]
public class StockPosition {
  private String symbol;
  private double quantity;
}

public class StockLot {
  private String iotNumber;
  private Order buyingOrder;

  public double getBuyingPrice();
}
\end{lstlisting}

\subsubsection{Order}\label{zQGXHbZZH18ujr5z8za9e}
\\
Members can place stock trading orders to sell or acquire \texttt{StockPosition}.

\textbf{OrderPart, Order, and its derived classes}

\begin{lstlisting}[language=Java]
public class OrderPart {
  private double price;
  private double quantity;
  private Date executedAt;
}

// Order is an abstract class
public abstract class Order {
  private String orderNumber;
  public boolean isBuyOrder;
  private OrderStatus status;
  private TimeEnforcementType timeEnforcement;
  private Date creationTime;
  private HashMap<int, OrderPart> parts;

  public void setStatus(OrderStatus status);
  public boolean saveInDatabase();
  public void addOrderParts(OrderParts parts);
}

public class LimitOrder extends Order {
  private double priceLimit;
}

public class StopLimitOrder extends Order {
  private double priceLimit;
}

public class StopLossOrder extends Order {
  private double priceLimit;
}

public class MarketOrder extends Order {
  private double priceLimit;
}
\end{lstlisting}

\subsubsection{Transfer money}\label{ZqojUdCwYjrr2cT6ZJ70w}
\\
Members should be able to deposit and withdraw money via \texttt{Check}, \texttt{Wire}, or \texttt{ElectronicBank} transfer.

\textbf{TransferMoney and its derived classes}

\begin{lstlisting}[language=Java]
// TransferMoney is an abstract class
public abstract class TransferMoney {
  private int id;
  private Date creationDate;
  public int fromAccount;
  private int toAccount;

  public abstract boolean initiateTransaction();
}

public class ElectronicBank extends TransferMoney {
  private String bankName;
  
  public boolean initiateTransaction();
}

public class Wire extends TransferMoney {
  private int wire;

  public boolean initiateTransaction();
}

public class Check extends TransferMoney {
  private String checkNumber;

  public boolean initiateTransaction();
}

public class DepositMoney {
  private int transactionId;

  public boolean initiateTransaction();
}

public class WithdrawMoney {
  private int transactionId;
  public boolean initiateTransaction();
}
\end{lstlisting}

\subsubsection{Notification}\label{JbRss2Szx199OeoiK9EUr}
\\
The \texttt{Notification} class is another abstract class responsible for sending notifications to the users, with the \texttt{SmsNotification} and \texttt{EmailNotification} classes as its child classes. The implementation of this class is shown below:

\textbf{Notification and its derived classes}

\begin{lstlisting}[language=Java]
public abstract class Notification {
    private String notificationId;
    private Date creationDate;
    private String content;

    public abstract boolean sendNotification();
}

public class SmsNotification extends Notification {
    private int phoneNumber;

    public boolean sendNotification();
}

public class EmailNotification extends Notification {
    private String email;

    public boolean sendNotification();
}
\end{lstlisting}

\subsubsection{Stock exchange}\label{z1ZmwTGKMg-zaIjVewlbk}
\\
The stock brokerage system will get all the stocks and their current pricing from the \texttt{StockExchange} class.

\textbf{The StockExchange class}

\begin{lstlisting}[language=Java]
public class StockExchange {
  // The StockExchange is a singleton class that ensures it will have only one active instance at a time
  private static StockExchange instance = null;
  private StockExchange() {}
  public static StockExchange getInstance();
  public boolean placeOrder(Order order);
  public boolean acknowledge(Order order);
}
\end{lstlisting}

\subsection{Executable code: Online stock brokerage system}\label{rrA_rpx8uURmutoS4W50Y}
\\
Below is a fully self-contained, runnable program in Java, C\#, C++, Python, and JavaScript that demonstrates the core workflows of the online stock brokerage system. The main driver code of the system resides in the~\texttt{Driver.java},~\texttt{Driver.cs},~\texttt{Driver.py},~\texttt{Driver.cpp}, and~\texttt{Driver.js}~for all respective languages. You can click the ``Run'' button to execute the codes.

\textbf{Online Stock Brokerage System}

\begin{lstlisting}[language=Java]
import java.util.*;

public class Driver {
    public static void main(String[] args) {
        System.out.println("========== ONLINE STOCK BROKERAGE SYSTEM ==========");

        // --------------------------------------------------
        // SYSTEM INITIALIZATION
        // --------------------------------------------------
        System.out.println("
--- SYSTEM INITIALIZATION ---");

        // Create and configure essential system components
        StockInventory inventory = new StockInventory();
        StockExchange exchange = StockExchange.getInstance();

        // Create and initialize member
        Member member = new Member();
        member.setName("Alice");
        member.setAvailableFundsForTrading(10000.0);

        System.out.println("[INFO] Member Initialized: " + member.getName());
        System.out.println("[INFO] Available Funds: $" + member.getAvailableFundsForTrading());

        // --------------------------------------------------
        // SCENARIO 1: SEARCHING AND SELECTING A STOCK
        // --------------------------------------------------
        System.out.println("
--- SCENARIO 1: SEARCHING AND SELECTING A STOCK ---");

        String symbolToSearch = "AAPL";
        System.out.println("[STEP] Searching for stock with symbol: " + symbolToSearch);
        Stock stock = inventory.searchSymbol(symbolToSearch);

        // Fallback if stock not found
        if (stock == null) {
            System.out.println("[WARNING] Stock not found in inventory. Creating mock stock data.");
            stock = new Stock();
            stock.setSymbol(symbolToSearch);
            stock.setPrice(180.50);
        }

        System.out.println("[RESULT] Stock Found: " + stock.getSymbol() + " @ $" + stock.getPrice());

        // Member selects the stock
        member.selectStock(stock.getSymbol());
        double quantity = 10;
        System.out.println("[ACTION] Selected Quantity: " + quantity);

        // --------------------------------------------------
        // SCENARIO 2: PLACING A LIMIT ORDER
        // --------------------------------------------------
        System.out.println("
--- SCENARIO 2: PLACING A LIMIT ORDER ---");

        Order order = new LimitOrder();
        order.setIsBuyOrder(true);
        order.setOrderNumber("ORD123");
        ((LimitOrder) order).setPriceLimit(stock.getPrice() - 5); // Buy at $175.50
        order.setTimeEnforcement(TimeEnforcementType.GOOD_TILL_CANCELED);
        order.setCreationTime(new Date());

        System.out.println("[ORDER] Type: Limit Order");
        System.out.println("[ORDER] Order Number: " + order.getOrderNumber());
        System.out.println("[ORDER] Buy Order: " + order.isBuyOrder);
        System.out.println("[ORDER] Limit Price: $" + ((LimitOrder) order).getPriceLimit());
        System.out.println("[ORDER] Time Enforcement: " + order.getTimeEnforcement());
        System.out.println("[ORDER] Created On: " + order.getCreationTime());

        // Place order through stock exchange
        boolean orderPlaced = exchange.placeOrder(order);
        inventory.sendOrderDetails(order);

        if (orderPlaced) {
            System.out.println("[SUCCESS] Order placed successfully: " + order.getOrderNumber());
        } else {
            System.out.println("[FAILED] Order placement failed.");
        }

        // Acknowledge order by exchange
        boolean ack = exchange.acknowledge(order);
        System.out.println("[EXCHANGE] Order Acknowledged: " + ack);

        // --------------------------------------------------
        // SCENARIO 3: SENDING ORDER NOTIFICATION
        // --------------------------------------------------
        System.out.println("
--- SCENARIO 3: SENDING ORDER NOTIFICATION ---");

        Notification notification = new EmailNotification();
        notification.setCreationDate(new Date());
        notification.setContent("Your order " + order.getOrderNumber() + " has been received.");
        ((EmailNotification) notification).setEmail("alice@example.com");

        boolean sent = notification.sendNotification();
        System.out.println("[NOTIFICATION] Email sent to " + ((EmailNotification) notification).getEmail() + ": " + sent);

        // --------------------------------------------------
        // SESSION END
        // --------------------------------------------------
        System.out.println("
===================================================");
        System.out.println("   Thank you for trading with us, " + member.getName());
        System.out.println("===================================================");
    }
}
\end{lstlisting}

\subsubsection{What does this code show}\label{J3bG_llM0NV0b5TuMrvue}

\begin{itemize}
\item
\protect\phantomsection\label{QNhNF0lWcoHL4LCsans-7}
\textbf{System initialization:}~Initializes the core components of the online stock brokerage system, including the stock inventory, stock exchange (as a Singleton instance), and a member named Alice with a starting trading fund of \$10,000. This sets up the environment for simulating stock trading operations.
\item
\protect\phantomsection\label{Ja9Pb1jptavE7CAZAYazm}
\textbf{Scenario 1 (Stock search and selection):}~Alice searches for a stock using its symbol (``AAPL''). If the stock is found in the inventory, its details are retrieved; otherwise, a mock stock is created. Alice selects this stock and specifies the quantity to trade, simulating a typical user action before placing an order.
\item
\protect\phantomsection\label{37Nusdx2j7vxmAjmHorJJ}
\textbf{Scenario 2 (Placing a limit order):}~Alice creates a buy limit order for the selected stock, specifying a limit price below the current market value and setting the order enforcement to ``Good Till Canceled.'' The order is submitted to the stock exchange, and its details are sent to the inventory. The system confirms whether the order was placed successfully and receives an acknowledgment from the exchange.
\item
\protect\phantomsection\label{Ta_ogv30GFWYKXifp8kp2}
\textbf{Scenario 3 (Notification of order placement):}~Once the order is placed and acknowledged, the system generates an email notification for Alice, confirming that her order has been received. This demonstrates how real-time notifications keep users informed of important transaction events.
\end{itemize}
\begin{quote}
\textbf{Hint: Try customizing the code!}\\
This code is designed for experimentation, not just reading. You can edit, extend, or modify any part of the example.
\end{quote}
\\
In this chapter, we\textquotesingle ve explored the complete design of an online stock brokerage system. We 've also looked at how a basic system can be visualized using various UML diagrams and designed using object-oriented principles and design patterns.

% End of section content